\section{Un tablier tiré des notes de Grothendieck: les pseudo-corps}\label{sec:pseudofields}

Les précédents du \S\ref{sec:precedents} montrent Grothendieck lisant des invariants d'un topos sur deux présentations. Ses notes d'algèbre universelle contiennent aussi un tablier que le registre ne contenait pas encore: le topos des \q{pseudo-corps}, dont il donne trois présentations.

\subsection{Le tablier}

Dans le dossier 161-2 \cite{grothendieck-161-2}, pp.~38--43, Grothendieck énumère des propriétés d'anneaux et, pour chacune, un sous-topos du topos $X_{\mathrm{ann}}\simeq\widehat R$ qui classifie les anneaux, $R$ désignant la catégorie des schémas affines de type fini sur $\mathbb Z$. Sa numérotation change d'une page à l'autre: à la p.~41, l'item~12 est celui des corps, tandis qu'à la p.~42 il renumérote, de sorte que l'item~12 devient celui des anneaux compacts et l'item~13 celui des anneaux compacts réduits. Dans la suite, \q{item~12} et \q{item~13} renvoient à la numérotation de la p.~42, et nous parlons des \q{corps de la p.~41}. L'item~13 est l'\emph{anneau compact réduit}, qu'il nomme \emph{pseudo-corps}: pour tout $f$, il existe un unique idempotent $e$ tel que $ef$ soit inversible dans $eA$ et $(1-e)f=0$, de sorte que $A\simeq A_f\times A/fA$. Ce sont les anneaux commutatifs réguliers au sens de von Neumann. Il donne trois présentations du sous-topos qu'ils définissent:
\[
X_{\mathrm{ps.\,corps}}\;\simeq\;\widetilde R_{\mathrm{ps.\,corps}}\;\simeq\;\widehat{R_{\mathrm{cons}}},
\]
les faisceaux sur $R$ pour la topologie où une famille est couvrante quand l'un de ses membres a une section fibre par fibre, et les préfaisceaux sur la catégorie $R_{\mathrm{cons}}$ obtenue à partir de $R$ en inversant les flèches bijectives à extensions résiduelles triviales (pp.~41--42).

\subsection{Vérification de la présentation par préfaisceaux}

La théorie des pseudo-corps s'axiomatise par les axiomes des anneaux commutatifs et le séquent
\[
\forall a\ \exists!\,b\ \ (a^2b = a\ \wedge\ ab^2 = b):
\]
dans un anneau commutatif, le quasi-inverse $b$ de $a$ est unique, et l'idempotent de la définition de Grothendieck est $e = ab$. La théorie est donc cartésienne, et par suite de type préfaisceau: son topos classifiant est le topos des foncteurs à valeurs dans $\Set$ sur ses modèles de présentation finie. Notons $T$ l'adjoint à gauche de l'inclusion des anneaux réguliers au sens de von Neumann dans les anneaux. Cette inclusion commute aux colimites filtrantes, donc $T$ préserve la présentation finie. Réciproquement, tout pseudo-corps de présentation finie est de la forme $T(A)$, avec $A$ un anneau de présentation finie: dans une présentation finie, on remplace chaque quasi-inverse par une nouvelle variable soumise à ses deux équations. D'après Olivier \cite{olivier}, $\operatorname{Spec}T(A)$ est $\operatorname{Spec}A$ muni de sa topologie constructible, et les anneaux locaux de $T(A)$ sont les corps résiduels de $A$. D'après Pierce \cite{pierce}, un anneau régulier au sens de von Neumann est l'anneau des sections globales du faisceau de ses anneaux locaux sur son spectre, qui est un espace booléen. Ainsi, pour $u:A\to B$ de présentation finie, $T(u)$ est un isomorphisme si et seulement si $\operatorname{Spec}u$ est bijectif et toutes les extensions résiduelles sont triviales: une bijection continue entre espaces booléens est un homéomorphisme, et un morphisme de faisceaux qui est un isomorphisme sur les fibres est un isomorphisme. Les flèches que $T$ inverse sont donc exactement celles de la classe $W$ qu'inverse Grothendieck.

Reste à voir que les pseudo-corps de présentation finie s'obtiennent à partir des anneaux de présentation finie en inversant $W$; voici notre argument. La construction d'Olivier écrit $T(B)$ comme colimite filtrante de $B$-algèbres de présentation finie obtenues en adjoignant des quasi-inverses, un élément $a$ à la fois:
\[
B' = B[y]/(a^2y - a,\ ay^2 - y).
\]
Chaque morphisme $B\to B'$ de cette tour est bijectif sur les spectres, à extensions résiduelles triviales: au-dessus d'un point où $a$ est non nul dans le corps résiduel, les équations imposent $y = 1/a$, et là où $a = 0$, elles imposent $y = 0$. Ces morphismes sont donc dans $W$, de même que leurs composés. Comme $A$ est de présentation finie, un morphisme $T(A)\to T(B)$, c'est-à-dire un morphisme $A\to T(B)$, se factorise par un étage de la tour; il est donc représenté par un toit $A\to B'\leftarrow B$, avec $B\to B'$ dans $W$. La catégorie des pseudo-corps de présentation finie est ainsi la localisée de la catégorie des anneaux de présentation finie par rapport à $W$; son opposée est $R_{\mathrm{cons}}$, et le topos classifiant est $\widehat{R_{\mathrm{cons}}}$, comme l'énonce Grothendieck. Nous n'avons pas vérifié la première présentation, par la topologie $\mathrm{ps.\,corps}$, et rien dans cette section n'est formalisé.

\subsection{Les arches sur ce tablier}

\begin{itemize}
\item \emph{De Morgan:} oui, car les modèles de présentation finie d'une théorie cartésienne ont des sommes amalgamées, et donc s'amalgament (arche de De Morgan, \S\ref{sec:recall}).
\item \emph{Deux valeurs:} non, car l'énoncé $1=0$ définit un sous-objet de l'objet final qui n'est ni $0$ ni $1$.
\item \emph{Connexe:} oui, car $T(\mathbb Z)$ est un modèle de présentation finie initial, de sorte que la catégorie de ces modèles est connexe.
\item \emph{Booléen:} non, car le morphisme $T(\mathbb Z)\to T(\mathbb F_2)$ n'est pas inversible.
\item \emph{Points essentiels:} les pseudo-corps de présentation finie $T(A)$, à isomorphisme près.
\item \emph{Sous-topos:} les théories quotients de la théorie des pseudo-corps, en vertu de la dualité entre théories quotients et sous-topos \cite{caramello-tst}. Plusieurs items de la liste de Grothendieck en font partie: les corps de la p.~41, qu'il identifie aux pseudo-corps intègres et aux pseudo-corps locaux (il marque la dernière égalité \q{à prouver}, et elle est vraie, puisqu'un anneau local régulier au sens de von Neumann est un corps); les corps séparablement clos et algébriquement clos (items~14--15); les anneaux annulés par $n$ et ceux où $n$ est inversible (item~16, p.~43), qu'il marque comme sous-topos ouverts.
\end{itemize}

\subsection{Réponse à une question de la marge}

Pour l'item~12, les anneaux \emph{compacts} (tout $f$ inversible sur un idempotent $e$ et nilpotent sur $1-e$, p.~42), Grothendieck note en marge que la topologie n'est pas cohérente et demande si le topos a assez de points. À la p.~41, il avertit aussi que les sous-topos de sa liste sont déterminés par leurs points, mais que ce n'est pas clair pour leurs intersections, de sorte que vérifier une inclusion entre intersections sur les structures ensemblistes pourrait ne pas suffire. Nous levons ces deux doutes en passant par les axiomes. Chaque item de sa liste s'axiomatise, sur la signature dénombrable des anneaux, par un nombre fini de séquents géométriques. Pour les anneaux compacts, l'un de ces séquents comporte une disjonction dénombrable, puisque $(1-e)f$ est nilpotent d'un certain ordre $n$, sans borne sur $n$; c'est pourquoi la topologie n'est pas cohérente. Le topos classifiant d'une théorie géométrique dénombrable est le topos des faisceaux sur son site syntactique, catégorie dénombrable à limites finies dont la topologie est engendrée par une famille dénombrable de cribles. Or un théorème de Makkai et Reyes \cite[th.~6.2.4]{makkai-reyes} (voir aussi \cite[cor.~3.3.10]{diliberti-rogers}) affirme que le topos des faisceaux sur une catégorie dénombrable à produits fibrés, pour une topologie engendrée par une famille dénombrable de cribles, a assez de points. Le topos des anneaux compacts a donc assez de points; celui des pseudo-corps en a trivialement assez, puisque c'est un topos de préfaisceaux. Dès qu'un sous-topos $E'$ a assez de points, une inclusion $E'\subseteq E''$ se vérifie sur les points. Une intersection d'une famille dénombrable de sous-topos de sa liste classifie la théorie obtenue en réunissant leurs axiomes, qui reste dénombrable; elle a donc assez de points, et les vérifications que Grothendieck fait sur les points sont légitimes pour elle. L'argument ne couvre pas les intersections d'une famille non dénombrable de sous-topos.

\subsection{Les sous-topos ouverts, décrits par des corps}\label{sec:pseudofields-opens}

Dans la présentation par préfaisceaux, un sous-objet de l'objet final de $X_{\mathrm{ps.\,corps}}$ est une classe $U$ de pseudo-corps de présentation finie stable vers le haut par homomorphismes: si $T(A)\in U$ et s'il existe un morphisme $T(A)\to T(B)$, alors $T(B)\in U$. Le cadre des ouverts est donc celui des parties stables vers le haut du préordre \q{il existe un morphisme $T(A)\to T(B)$}, et le décrire revient à décider quand un tel morphisme existe. L'item~16 de Grothendieck (p.~43) marque deux de ces ouverts, les anneaux annulés par $n$ et ceux où $n$ est inversible. Nous décidons l'existence des morphismes par les corps résiduels, et en tirons une description de tous les ouverts. Les anneaux sont commutatifs et unitaires, \q{régulier} veut dire régulier au sens de von Neumann, et pour un anneau $A$ on pose
\[
F(A) = \{\text{corps } K \text{ tels qu'il existe un morphisme d'anneaux } A\to K\}.
\]
L'appartenance de $K$ à $F(A)$ ne dépend que des sous-corps de $K$ de type fini sur son corps premier (l'image de $A$ en engendre un quand $A$ est de type fini), de sorte que, pour $A$ de type fini, on peut voir $F(A)$ comme un ensemble de classes d'isomorphisme de corps de type fini, sans changer les inclusions entre les $F(A)$.

\paragraph{Rappels sur un anneau régulier $R$.} Nous n'utilisons que ce qui suit, élémentaire. Pour un idempotent $e$, $Re$ est un anneau d'unité $e$, et $r\mapsto er$ est un morphisme d'anneaux $R\to Re$; les parties $D(e) = \{x\in\operatorname{Spec}R : e\notin x\}$ sont ouvertes et fermées, et $\operatorname{Spec}R$ est quasi-compact. Pour un idéal premier $x$, l'anneau local $R_x$ est le corps résiduel $R/x$: si $s\notin x$ et si $s^*$ est le quasi-inverse de $s$, alors $e = ss^*$ est un idempotent tel que $s = se$, donc $e\notin x$, d'où $1-e\in x$; par suite $r/s = rs^*/1$ dans $R_x$, et $r/1 = 0$ dans $R_x$ si et seulement si $er = 0$ pour un idempotent $e\notin x$ (si $sr = 0$, prendre $e = ss^*$). Enfin, un idempotent contenu dans tous les idéaux premiers est nilpotent, donc nul.

\begin{lemma}[du local au global]\label{lem:pf-local}
Soit $A = \mathbb Z[x_1,\dots,x_n]/(f_1,\dots,f_m)$ un anneau de présentation finie et $R$ un anneau régulier (pas nécessairement de présentation finie). Si chaque corps résiduel $R/x$, $x\in\operatorname{Spec}R$, reçoit un morphisme d'anneaux depuis $A$, alors $R$ en reçoit un. Plus précisément, pour tout $x$, il existe un idempotent $e_x\notin x$ et un morphisme d'anneaux $A\to Re_x$.
\end{lemma}

\begin{proof}
Fixons $x$ et un morphisme $A\to R/x$, c'est-à-dire $a\in(R/x)^n$ tel que $f_j(a) = 0$ pour tout $j$. Comme $R\to R_x = R/x$ est surjectif, relevons $a$ en $r\in R^n$. Chaque $f_j(r)$ a une image nulle dans $R_x$, donc $e_{x,j}f_j(r) = 0$ pour un idempotent $e_{x,j}\notin x$; le produit $e_x$ des $m$ idempotents $e_{x,j}$ est un idempotent qui n'est pas dans $x$, puisque $x$ est premier. Le morphisme $\mathbb Z[x_1,\dots,x_n]\to Re_x$ qui envoie $x_i$ sur $e_xr_i$ est le composé de $x_i\mapsto r_i$ et de $R\to Re_x$; il envoie donc $f_j$ sur $e_xf_j(r) = 0$, et se factorise par $A$. C'est l'étape locale, et elle utilise les deux conditions de finitude sur $A$: un nombre fini de générateurs pour relever, un nombre fini de relations pour trouver un seul idempotent.

Les ouverts $D(e_x)$ recouvrent $\operatorname{Spec}R$, qui est quasi-compact; un nombre fini suffit, soit $D(e_1),\dots,D(e_k)$, avec des morphismes $\varphi_i\colon A\to Re_i$. Posons $\varepsilon_1 = e_1$ et $\varepsilon_i = e_i(1-e_1)\cdots(1-e_{i-1})$. Ces idempotents sont deux à deux orthogonaux, $\varepsilon_ie_i = \varepsilon_i$, et $1-\sum_i\varepsilon_i = \prod_i(1-e_i)$ est un idempotent contenu dans tout idéal premier (chaque premier évite un $e_i$, donc contient $1-e_i$): il est nul. Ainsi $R\cong\prod_iR\varepsilon_i$, et les composés $A\xrightarrow{\varphi_i}Re_i\xrightarrow{\ \cdot\varepsilon_i\ }R\varepsilon_i$ se recollent en un morphisme d'anneaux $A\to R$.
\end{proof}

On ne peut pas supprimer la finitude de $A$. Soit $P_1\sqcup P_2$ une partition des nombres premiers en deux parties infinies, et soit $A$ l'anneau engendré par $e$, des $y_q$ ($q\in P_1$) et des $z_q$ ($q\in P_2$), soumis à $e^2 = e$, $1-e = qy_q$ pour $q\in P_1$ et $e = qz_q$ pour $q\in P_2$. Les corps résiduels de $R = T(\mathbb Z)$ sont $\mathbb Q$ et les $\mathbb F_q$, et chacun reçoit un morphisme depuis $A$ ($e = 1$ en $q\in P_1$, $e = 0$ en $q\in P_2$, n'importe quoi sur $\mathbb Q$). Mais un morphisme $A\to T(\mathbb Z)$ donnerait un idempotent de $T(\mathbb Z)$ égal à $1$ exactement aux premiers de $P_1$, alors que les idempotents de $T(\mathbb Z)$ sont les parties ouvertes et fermées de $\operatorname{Spec}\mathbb Z$ muni de la topologie constructible \cite{olivier}, c'est-à-dire les ensembles finis de premiers et leurs complémentaires.

\begin{theorem}[critère]\label{th:pf-criterion}
Soient $A$ un anneau de présentation finie et $B$ un anneau quelconque. Les conditions suivantes sont équivalentes:
\begin{enumerate}
\item il existe un morphisme d'anneaux $T(A)\to T(B)$;
\item chaque corps résiduel $\kappa(\mathfrak p)$, $\mathfrak p\in\operatorname{Spec}B$, reçoit un morphisme d'anneaux depuis $A$;
\item $F(B)\subseteq F(A)$: tout corps qui reçoit un morphisme depuis $B$ en reçoit un depuis $A$.
\end{enumerate}
\end{theorem}

\begin{proof}
Par adjonction, (1) dit qu'il existe un morphisme $A\to T(B)$. D'après Olivier \cite{olivier}, $\operatorname{Spec}T(B)\to\operatorname{Spec}B$ est bijectif et le corps résiduel de $T(B)$ au point au-dessus de $\mathfrak p$ est $\kappa(\mathfrak p)$. Donc (1)$\Rightarrow$(2) en composant avec $T(B)\to\kappa(\mathfrak p)$, et (2)$\Rightarrow$(1) est le lemme~\ref{lem:pf-local} pour $R = T(B)$. Pour (2)$\Leftrightarrow$(3): les corps résiduels de $B$ sont dans $F(B)$, un morphisme $B\to K$ vers un corps se factorise par le corps résiduel en son noyau, et $F(A)$ est stable par extension de corps.
\end{proof}

La nécessité dans~(2) est immédiate; ce qu'apporte le lemme~\ref{lem:pf-local} est la suffisance, et les corps résiduels n'ont besoin d'être ni finis, ni en nombre fini. L'anneau $B$ n'a pas besoin d'être de présentation finie, et le critère dit que $T(A)\to T(B)$ existe si et seulement si $T(B)$ satisfait l'énoncé primitif positif $\sigma_A = \exists\bar x\,(f_1(\bar x) = \dots = f_m(\bar x) = 0)$, si et seulement si chaque corps résiduel de $B$ le satisfait.

\paragraph{Le treillis $D$.} Soit $D$ l'ensemble des classes $F(A)$, $A$ de présentation finie. On a $F(0) = \emptyset$, $F(\mathbb Z)$ est la classe de tous les corps, $F(A)\cap F(A') = F(A\otimes A')$ et $F(A)\cup F(A') = F(A\times A')$, la dernière égalité parce qu'un corps n'a pas d'autre idempotent que $0$ et $1$. Ainsi $D$ est un treillis distributif d'ensembles de corps. C'est le treillis de Lindenbaum des énoncés cohérents de la théorie des corps: modulo cette théorie, tout énoncé cohérent est une disjonction finie d'énoncés $\sigma_A$, et $\sigma_A\vee\sigma_{A'}$ équivaut à $\sigma_{A\times A'}$ dans un corps. Dans un pseudo-corps, $\sigma_A\vee\sigma_{A'}$ et $\sigma_{A\times A'}$ ne sont plus équivalents: $\mathbb F_2\times\mathbb F_3$ satisfait $\sigma_{\mathbb Z/6}$, mais ni $\sigma_{\mathbb Z/2}$, ni $\sigma_{\mathbb Z/3}$. C'est toute la différence entre les deux cadres ci-dessous.

\begin{theorem}[le cadre des ouverts]\label{th:pf-frame}
\begin{enumerate}
\item L'application $T(A)\mapsto F(A)$ identifie la réflexion ordonnée de la catégorie des pseudo-corps de présentation finie à l'opposé de $D$: il existe un morphisme $T(A)\to T(B)$ si et seulement si $F(B)\subseteq F(A)$.
\item Le cadre des sous-objets de l'objet final de $X_{\mathrm{ps.\,corps}}$ est isomorphe au cadre $\mathrm{Down}(D)$ de toutes les parties de $D$ stables vers le bas, ordonnées par inclusion. À une partie $S$ stable vers le bas correspond l'ouvert $U_S$, vrai dans un pseudo-corps $R$ (de présentation finie ou non) si et seulement si un membre de $S$ contient tous les corps résiduels de $R$; le sous-topos ouvert $o(U_S)$ classifie ces pseudo-corps.
\item La partie principale ${\downarrow}F(A)$ correspond à $U_A$, image du foncteur représenté par $T(A)$, qui est l'ouvert défini par $\sigma_A$. Une partie $S$ stable vers le bas est un idéal (stable par réunions finies) si et seulement si $U_S$ est déterminé par les corps résiduels, c'est-à-dire: $R\in U_S$ si et seulement si chaque corps résiduel de $R$ est dans $U_S$.
\item Le sous-topos $X_{\mathrm{corps}}$ des corps (p.~41) a pour cadre des ouverts $\mathrm{Idl}(D)$, et la restriction des ouverts de $X_{\mathrm{ps.\,corps}}$ à $X_{\mathrm{corps}}$ envoie $S$ sur l'idéal qu'elle engendre.
\end{enumerate}
\end{theorem}

\begin{proof}
(1) est le théorème~\ref{th:pf-criterion}, et chaque membre de $D$ est un $F(A)$. (2) Dans un topos de préfaisceaux $[\mathcal C,\Set]$, les sous-objets de l'objet final sont les classes d'objets de $\mathcal C$ stables vers le haut, donc, d'après~(1), les parties de $D$ stables vers le bas: on pose $S_U = \{F(B) : T(B)\in U\}$. Tout pseudo-corps $R$ est colimite filtrante de pseudo-corps de présentation finie, de sorte que $R\in U$ si et seulement si un $T(B)\in U$ s'envoie dans $R$, si et seulement si (lemme~\ref{lem:pf-local} et nécessité du théorème~\ref{th:pf-criterion}, valables pour tout $R$ régulier) tous les corps résiduels de $R$ sont dans $F(B)$ pour un $F(B)\in S_U$. (3) Le premier énoncé est~(1). Soit $S$ un idéal, et supposons chaque corps résiduel $R/x$ contenu dans un membre $F(A_x)$ de $S$. Le lemme~\ref{lem:pf-local} donne des idempotents $e_x\notin x$ et des morphismes $A_x\to Re_x$; comme dans sa preuve, un nombre fini recouvre, on les orthogonalise en $\varepsilon_1,\dots,\varepsilon_k$, et $R\cong\prod R\varepsilon_i$ reçoit un morphisme depuis $\prod A_{x_i}$, dont la classe $\bigcup F(A_{x_i})$ est dans $S$. Réciproquement, si $U_S$ est déterminé par les corps résiduels et si $F(A), F(A')\in S$, les corps résiduels de $T(A\times A')$ sont ceux de $A$ et de $A'$, tous dans $U_S$, donc $F(A)\cup F(A') = F(A\times A')\in S$. (4) $X_{\mathrm{corps}}$ classifie la théorie cohérente des corps; c'est donc le topos des faisceaux sur sa catégorie syntactique cohérente pour la topologie cohérente, et un sous-objet de l'objet final y est un crible fermé sur $1$, c'est-à-dire un ensemble d'énoncés stable par conséquence et par disjonctions finies: un idéal du treillis de Lindenbaum $D$. La restriction de $U_S$ aux corps est l'énoncé $\bigvee_{C\in S}C$, dont l'idéal est engendré par $S$.
\end{proof}

D'après~(4), les idéaux de $D$ sont les ouverts de l'espace spectral $\operatorname{Spec}D$ des filtres premiers, et chaque filtre premier est l'ensemble des membres de $D$ qui contiennent un corps donné (par le théorème de compacité: pour un filtre premier $\mathfrak p$, les énoncés $\sigma_A$, $F(A)\in\mathfrak p$, et $\neg\sigma_B$, $F(B)\notin\mathfrak p$, sont finiment satisfaisables dans les corps, puisque $F(A_1\otimes\dots\otimes A_k)\not\subseteq F(B_1\times\dots\times B_l)$). Ainsi $\operatorname{Spec}D$ est l'espace des types existentiels positifs de corps, et le cadre de $X_{\mathrm{ps.\,corps}}$ est le cadre de toutes les parties stables vers le bas de son treillis d'ouverts quasi-compacts. En termes du faisceau de Pierce \cite{pierce}: $R\in U_S$ si et seulement si les fibres du faisceau de Pierce de $R$, au-dessus de son spectre booléen, sont toutes dans un même membre de $S$. Tout ouvert non vide contient l'anneau nul, de sorte que deux ouverts non vides se rencontrent: c'est la propriété de De Morgan du topos de préfaisceaux, vue sur les ouverts.

\paragraph{Petits cas.} (a) \emph{Caractéristiques.} Envoyer $F(A)$ sur l'ensemble des caractéristiques de ses corps, c'est-à-dire sur l'image de $\operatorname{Spec}A\to\operatorname{Spec}\mathbb Z$, est un morphisme de treillis surjectif de $D$ sur le treillis des parties constructibles de $\operatorname{Spec}\mathbb Z$ (ensembles finis de premiers, et leurs complémentaires augmentés de~$0$); l'image d'un produit fibré est l'intersection des images, et la classe des corps de caractéristique dans une partie constructible $C$ est $F(\mathbb Z[1/m]\times\mathbb Z/n)$ pour des entiers sans facteur carré $m$, $n$ convenables. C'est une section, de sorte que les ouverts qui ne dépendent que de la caractéristique forment le cadre des parties stables vers le bas du treillis des constructibles de $\operatorname{Spec}\mathbb Z$. Le morphisme n'est pas injectif: pour $A = \mathbb Z[x]/((x^2-2)(x^2-3)(x^2-6))$, tout corps de caractéristique positive est dans $F(A)$, $\mathbb Q$ n'y est pas, et $\mathbb Q(\sqrt2)$ y est.

(b) \emph{L'item 16.} Dans $X_{\mathrm{ps.\,corps}}$, \q{$n = 0$} est $U_{\mathbb Z/n} = {\downarrow}V(n)$ et \q{$n$ inversible} est $U_{\mathbb Z[1/n]} = {\downarrow}D(n)$, en notant $V(n)$, $D(n)$ les classes de corps de caractéristique divisant, resp.\ ne divisant pas,~$n$. Le premier ne dépend que du radical de $n$ ($4 = 0$ et $2 = 0$ définissent le même ouvert, un pseudo-corps étant réduit), ce qui n'est pas le cas dans $\widehat R$. Leur intersection est $U_0$, le point ouvert de l'anneau nul; pour $n = 1$, c'est le N.B.\ de Grothendieck, $X_{1\text{-tors}} = X_\emptyset$, topos ponctuel. Leur réunion ${\downarrow}V(n)\cup{\downarrow}D(n)$ n'est pas le topos entier: pour $n = 2$, $\mathbb F_2\times\mathbb Q$ n'est dans aucun des deux. Mais l'idéal qu'ils engendrent est $D$ tout entier, car $V(n)\cup D(n) = F(\mathbb Z/n\times\mathbb Z[1/n])$ est l'élément maximal; en effet tout pseudo-corps se scinde en $R/nR\times R[1/n]$, le long de l'idempotent $nn^*$. Ainsi, dans $X_{\mathrm{corps}}$, les deux ouverts sont complémentaires; dans $X_{\mathrm{ps.\,corps}}$, ils sont disjoints au point près et ne recouvrent pas, et seule leur clôture par produits finis est le tout. La partie de caractéristique nulle $\bigcap_nU_{\mathbb Z[1/n]}$ a pour intérieur $U_0$ (un membre de $D$ dont tous les corps sont de caractéristique nulle a un ensemble constructible, donc vide, de caractéristiques); ainsi, comme dans $\widehat R$ (p.~47), le sous-topos des pseudo-corps sur $\mathbb Q$ n'est pas ouvert.

(c) \emph{Sur $\mathbb F_p$.} Le sous-topos ouvert $o(U_{\mathbb F_p})$ a pour cadre $\mathrm{Down}(D_p)$, $D_p = {\downarrow}F(\mathbb F_p)$, formé des classes $F(A)$ pour $A$ de présentation finie sur $\mathbb F_p$. Sur les algèbres finies étales $B = \prod_i\mathbb F_{p^{d_i}}$, qui sont leur propre $T$, il existe un morphisme $B\to B'$ si et seulement si chaque $d'_j$ est multiple d'un $d_i$; ces points forment donc le treillis des parties de $(\mathbb N_{\ge1},\mid)$ stables vers le haut et finiment engendrées, ordonnées par l'inclusion renversée. $D_p$ n'est pas déterminé par les corps algébriques: $A = \mathbb F_p[x,1/(x^p-x)]$ s'envoie dans $\mathbb F_p(t)$ et dans chaque $\mathbb F_{p^d}$, $d\ge2$, mais pas dans $\mathbb F_p$.

(d) \emph{Sur $\mathbb Q$.} Le même théorème vaut pour les $\mathbb Q$-algèbres (le lemme~\ref{lem:pf-local} est inchangé), pour le topos des pseudo-corps sur $\mathbb Q$, préfaisceaux sur les $T(A)$, $A$ de présentation finie sur $\mathbb Q$. Sur les algèbres finies étales $B = \prod K_i$, $B' = \prod K'_j$, il existe un morphisme $B\to B'$ si et seulement si chaque $K'_j$ contient un conjugué d'un $K_i$; par exemple $\mathbb Q(i)$ et $\mathbb Q(\sqrt2)$ sont incomparables, et $\mathbb Q(i)\times\mathbb Q(\sqrt2)$ s'envoie dans l'un et l'autre.

\paragraph{Ce qui est démontré, et ce qui ne l'est pas.} Le lemme~\ref{lem:pf-local} et les théorèmes~\ref{th:pf-criterion} et~\ref{th:pf-frame} sont démontrés ci-dessus à partir de la description de $\operatorname{Spec}T(B)$ par Olivier et des faits standard sur les topos de préfaisceaux et les théories cohérentes rappelés dans les preuves. Ils sont de l'édition, non de Grothendieck: les pages énoncent la présentation $\widehat{R_{\mathrm{cons}}}$ et énumèrent des ouverts, sans décrire le cadre.

\leanentry{} Le lemme~\ref{lem:pf-local} est démontré dans \href{https://github.com/Commutator-IO/GrothendieckArchivesProofs/blob/main/Grothendieck/RingTheory/Folder161_2Pseudofield.lean}{\nolinkurl{Grothendieck/RingTheory/Folder161\_2Pseudofield.lean}} (GrothendieckArchivesProofs), pour tout anneau $A$ de présentation finie et tout anneau régulier $R$, avec les corps résiduels aux idéaux maximaux, qui sont ici tous les premiers: \lean{Grothendieck.Folder161_2.exists_ringHom_of_forall_isMaximal}, l'énoncé étant \lean{OpenProblems.Findings.Folder161_2.RingHomOfResidueFields}. Le pas local est \lean{Grothendieck.Folder161_2.exists_fg_le_of_ringHom}: un morphisme $A\to R/x$ se relève en $A\to R/K$, où $K\subseteq x$ est un idéal de type fini, égal à $(g)$ pour un idempotent $g$, d'où $e_x = 1-g$. La quasi-compacité devient la finitude d'une famille des $1-g$ engendrant l'idéal unité, et l'orthogonalisation se fait un idempotent à la fois, par le théorème chinois pour les idéaux comaximaux $(g)$ et $J+(1-g)$. Le contre-exemple, le théorème~\ref{th:pf-criterion} et le théorème~\ref{th:pf-frame} ne sont pas formalisés.

\paragraph{Est-ce connu?} Pas dans ce que nous avons lu. Caramello et Johnstone \cite[p.~2 et lemmes~1.1--1.2]{caramello-johnstone} définissent, pour un anneau régulier de présentation finie $R$ (de présentation finie dans la signature munie du quasi-inverse, ce qui revient, par l'argument ci-dessus, à $R = T(A)$ avec $A$ de présentation finie), l'ensemble $\operatorname{Char}R$ des caractéristiques des corps qui reçoivent un morphisme depuis $R$, montrent qu'il est fini sans $0$ ou cofini avec $0$, et se servent des sous-objets $T_A$ de l'objet final du topos classifiant des corps définis par \q{$\operatorname{Char}\subseteq A$} pour montrer que ce topos n'est pas de De Morgan. C'est l'image de $D$ dans les constructibles de $\operatorname{Spec}\mathbb Z$ du petit cas~(a), restreinte au topos des corps. Les \emph{Universal models and definability} de Caramello \cite[\S7(ii)]{caramello-universal} reprennent ce résultat pour les caractéristiques, et ses \emph{Lattices of theories} \cite[\S7.1]{caramello-lattices} donnent le dictionnaire général qu'utilise la preuve du théorème~\ref{th:pf-frame}: les sous-objets de l'objet final d'un topos de préfaisceaux sont les sous-catégories pleines stables par cribles, et ceux d'un topos classifiant sont les énoncés géométriques modulo la théorie. Aucun des trois ne décrit le cadre de l'un ou l'autre topos. Le lemme~\ref{lem:pf-local} est l'argument de recollement des produits booléens: Burris et Sankappanavar \cite[déf.~IV.8.1 et l'exemple de Dauns et Hofmann après le cor.~IV.8.14]{burris-sankappanavar} définissent la propriété de recollement (\emph{patchwork property}) et représentent les anneaux biréguliers, donc les anneaux commutatifs réguliers, comme produits booléens d'anneaux simples, c'est-à-dire de corps, et d'un anneau nul; mais ils n'énoncent pas qu'un énoncé primitif positif est vrai dans un produit booléen si et seulement s'il l'est dans chaque facteur. Cet énoncé pourrait bien se trouver chez Burris et Werner (Trans.\ Amer.\ Math.\ Soc.\ 248, 1979) ou chez Macintyre (Fund.\ Math.\ 81, 1973), que nous n'avons pas pu lire. Nous avons interrogé TheoremSearch le 23 novembre 2026, par neuf requêtes sur les énoncés primitifs positifs dans les produits booléens et les anneaux réguliers, et sur les sous-objets de l'objet final des topos classifiants des anneaux réguliers et des corps; aucun résultat n'énonce le lemme ni la description du cadre. Le \emph{Stone Spaces} et l'article \emph{Rings, fields and spectra} de Johnstone n'ont pas été consultés.
